%% datos.tex — Datos de tu trabajo. Escribe entre las llaves { }.
%% Los campos marcados como (opcional) pueden quedar vacíos, p. ej. \cotutor{},
%% o borrarse sin que nada falle. Los datos de la universidad (nombre, logo,
%% titulación, márgenes…) están en estilo/institucion.tex.
%% Desde Estudio TFG también puedes editarlos en el panel «Datos del trabajo».

\titulo{Título de tu trabajo}
\autor{Nombre Apellido1 Apellido2}
\email{nombre.apellido@universidad.es}          % (opcional)
\tutor{Dr. Nombre Apellido}
\cotutor{}                                      % (opcional)
\departamento{}                                 % (opcional) sale como «Departamento de …»
\intensificacion{}                              % (opcional) tu especialidad, mención o itinerario
\fecha{junio}{2026}                             % (opcional) mes y año de entrega; si falta, la fecha actual
\ciudad{}                                       % (opcional) sale en la página de créditos
\palabrasClave{palabra clave 1, palabra clave 2, palabra clave 3}  % (opcional) para el Resumen
\keywords{keyword 1, keyword 2, keyword 3}      % (opcional) las mismas, en inglés, para el Abstract

%% Ajustes del documento
\idioma{espanol}              % espanol | ingles (títulos automáticos: «Capítulo», «Índice»…)
\formato{impresion}           % impresion (doble cara, enlaces negros) | pantalla (una cara, enlaces en color)
\modo{borrador}               % borrador (marca «BORRADOR» al pie y notas \todo en rojo) | final (un \todo da error)
\estiloBibliografia{ieee}     % ieee (numérico, orden de cita) | apa (autor-año) | numeric | authoryear | alphabetic
\licencia{reservados}         % reservados | cc-by | cc-by-sa | cc-by-nc-sa | ninguna (texto de la página de créditos)
\atribucion{si}               % si | no (página final que cita la clase original de ARCO)
